% Generato da tools/notion2latex.py a partire dalla pagina Notion "lezione-4".

\chapter{Lezione 4 -- Sicurezza computazionale e OWF}

\emph{Segue gli appunti del prof. Le parti aggiunte per capire meglio sono segnate ``non negli appunti del prof''. PRG e indistinguibilità computazionale sono nella} Lezione 5

\section{Fine lezione 3: dimostrazione del lemma sulla collision probability}

\begin{inbreve}

Nella lezione 3, per dimostrare il Leftover Hash Lemma, avevamo usato un lemma senza dimostrarlo: \textbf{se una distribuzione collide quasi quanto l'uniforme, allora è quasi uniforme}. Il prof lo dimostra qui, usando la disuguaglianza di Cauchy-Schwarz.

\end{inbreve}

\subsection{Promemoria: cosa dice il lemma}

Sia \(Y\) una variabile aleatoria che prende valori in un insieme \(\mathcal{Y}\), e sia \(U\) uniforme su \(\mathcal{Y}\). Se

\[
\text{Col}(Y) \le \frac{1}{|\mathcal{Y}|}\,(1 + 4\varepsilon^2)
\]

allora \(\text{SD}(Y; U) \le \varepsilon\).

\begin{nota}

\textbf{Cosa sono le cose}

\begin{itemize}
\tightlist
\item
  \(\mathcal{Y}\) = l'insieme dei valori possibili; \(|\mathcal{Y}|\) = quanti sono.
\item
  \(U\) = la distribuzione uniforme su \(\mathcal{Y}\): ogni valore ha probabilità \(1/|\mathcal{Y}|\).
\item
  \(\text{Col}(Y) = \sum_{y} \Pr[Y = y]^2\) = probabilità di collisione: la probabilità che due estrazioni indipendenti di \(Y\) diano lo stesso valore. L'uniforme ha la più bassa possibile, \(1/|\mathcal{Y}|\).
\item
  \(\text{SD}(Y; U)\) = distanza statistica tra \(Y\) e l'uniforme (lezione 3).
\end{itemize}

\end{nota}

\subsection{Strumento: prodotto scalare e Cauchy-Schwarz}

\emph{(Definizioni non negli appunti del prof, servono per seguire la dimostrazione.)}

\textbf{Prodotto scalare} di due vettori \(\vec{u} = (u_1, \dots, u_N)\) e \(\vec{v} = (v_1, \dots, v_N)\): si moltiplicano le componenti a coppie e si somma.

\[
\langle \vec{u}, \vec{v} \rangle = u_1 v_1 + u_2 v_2 + \dots + u_N v_N
\]

\begin{esempio}

\textbf{Esempio}: \(\langle (1, 2), (3, 4) \rangle = 1 \cdot 3 + 2 \cdot 4 = 11\). In particolare \(\langle \vec{u}, \vec{u} \rangle = u_1^2 + \dots + u_N^2\).

\end{esempio}

\textbf{Disuguaglianza di Cauchy-Schwarz}: il prodotto scalare di due vettori non supera mai la radice del prodotto dei loro ``quadrati'':

\[
\langle \vec{u}, \vec{v} \rangle \le \sqrt{\langle \vec{u}, \vec{u} \rangle \cdot \langle \vec{v}, \vec{v} \rangle}
\]

\begin{esempio}

\textbf{Esempio}: con \((1, 2)\) e \((3, 4)\): a sinistra 11; a destra \(\sqrt{5 \cdot 25} = \sqrt{125} \approx 11{,}18\). Infatti \(11 \le 11{,}18\).

\end{esempio}

\subsection{Dimostrazione (del prof)}

\textbf{Passo 1: riscrivere la SD come prodotto scalare.} Per definizione:

\[
\text{SD}(Y; U) = \frac{1}{2} \sum_{y \in \mathcal{Y}} \Big| \Pr[Y = y] - \frac{1}{|\mathcal{Y}|} \Big|
\]

Il prof introduce due abbreviazioni, per ogni valore \(y\):

\begin{itemize}
\tightlist
\item
  \(q_y = \Pr[Y = y] - \frac{1}{|\mathcal{Y}|}\): di quanto la probabilità di \(y\) si scosta da quella uniforme (può essere positivo o negativo)
\item
  \(s_y = +1\) se \(q_y \ge 0\), \(s_y = -1\) altrimenti: il \textbf{segno} di \(q_y\)
\end{itemize}

Siccome un numero moltiplicato per il suo segno dà il suo valore assoluto (\(q_y \cdot s_y = |q_y|\)):

\[
\text{SD}(Y; U) = \frac{1}{2} \sum_{y \in \mathcal{Y}} q_y\, s_y = \frac{1}{2} \langle \vec{q}, \vec{s} \rangle
\]

dove \(\vec{q} = (q_y)_{y \in \mathcal{Y}}\) e \(\vec{s} = (s_y)_{y \in \mathcal{Y}}\) sono i vettori con una componente per ogni valore \(y\).

\textbf{Passo 2: Cauchy-Schwarz.}

\[
\text{SD}(Y; U) = \frac{1}{2} \langle \vec{q}, \vec{s} \rangle \le \frac{1}{2} \sqrt{\langle \vec{q}, \vec{q} \rangle \cdot \langle \vec{s}, \vec{s} \rangle} = \frac{1}{2} \sqrt{\sum_{y \in \mathcal{Y}} q_y^2 \cdot |\mathcal{Y}|}
\]

Infatti \(\langle \vec{s}, \vec{s} \rangle = \sum_y s_y^2 = |\mathcal{Y}|\), perché ogni \(s_y\) vale \(\pm 1\), quindi \(s_y^2 = 1\), e ce n'è uno per ogni valore.

\textbf{Passo 3: calcolare} \(\sum_y q_y^2\). Si sviluppa il quadrato \((a - b)^2 = a^2 + b^2 - 2ab\):

\[
\sum_{y} q_y^2 = \sum_{y} \Big( \Pr[Y = y] - \frac{1}{|\mathcal{Y}|} \Big)^2 = \sum_{y} \Big( \Pr[Y = y]^2 + \frac{1}{|\mathcal{Y}|^2} - 2\,\frac{\Pr[Y = y]}{|\mathcal{Y}|} \Big)
\]

Si sommano i tre pezzi separatamente:

\begin{itemize}
\tightlist
\item
  \(\sum_y \Pr[Y = y]^2 = \text{Col}(Y)\), per definizione
\item
  \(\sum_y \frac{1}{|\mathcal{Y}|^2} = |\mathcal{Y}| \cdot \frac{1}{|\mathcal{Y}|^2} = \frac{1}{|\mathcal{Y}|}\), perché è lo stesso numero sommato \(|\mathcal{Y}|\) volte
\item
  \(\sum_y 2\,\frac{\Pr[Y = y]}{|\mathcal{Y}|} = \frac{2}{|\mathcal{Y}|} \sum_y \Pr[Y = y] = \frac{2}{|\mathcal{Y}|}\), perché le probabilità sommano a 1
\end{itemize}

Quindi, usando l'ipotesi su \(\text{Col}(Y)\):

\[
\sum_{y} q_y^2 = \text{Col}(Y) + \frac{1}{|\mathcal{Y}|} - \frac{2}{|\mathcal{Y}|} \le \frac{1 + 4\varepsilon^2}{|\mathcal{Y}|} - \frac{1}{|\mathcal{Y}|} = \frac{4\varepsilon^2}{|\mathcal{Y}|}
\]

\textbf{Passo 4: rimettere insieme.}

\[
\text{SD}(Y; U) \le \frac{1}{2} \sqrt{\frac{4\varepsilon^2}{|\mathcal{Y}|} \cdot |\mathcal{Y}|} = \frac{1}{2} \sqrt{4\varepsilon^2} = \frac{1}{2} \cdot 2\varepsilon = \varepsilon
\]

∎

\textbf{Intuizione}: \(\sum_y q_y^2\) misura quanto \(Y\) si scosta dall'uniforme ``al quadrato'', ed è esattamente l'\textbf{eccesso di collisioni} \(\text{Col}(Y) - 1/|\mathcal{Y}|\). Se l'eccesso è piccolo, Cauchy-Schwarz garantisce che anche la distanza statistica è piccola.

\section{Sicurezza computazionale}

\begin{inbreve}

Finora la sicurezza valeva contro avversari con potenza di calcolo \textbf{infinita}, e il prezzo erano limiti pesantissimi (chiavi lunghe quanto i messaggi, chiavi che crescono con l'uso). Per avere crittografia \textbf{dimostrabilmente sicura ed efficiente} bisogna cambiare qualcosa: si fanno \textbf{due ipotesi} in più.

\end{inbreve}

I limiti visti finora:

\begin{itemize}
\tightlist
\item
  \textbf{Cifratura (Shannon)}: la chiave deve essere lunga almeno quanto il messaggio e va usata una volta sola.
\item
  \textbf{MAC}: per autenticare \(t\) messaggi servono chiavi di \((t+1)\lambda\) bit.
\item
  \textbf{Casualità}: da una fonte con min-entropia \(k\) non si estraggono più di circa \(k\) bit puliti.
\end{itemize}

Le due ipotesi del prof:

\begin{enumerate}
\def\labelenumi{\arabic{enumi}.}
\tightlist
\item
  \textbf{L'avversario gira su una macchina ``efficiente''}: non ha potenza infinita.
\item
  \textbf{Esistono problemi che sembrano ``difficili''}, cioè che nessuna macchina efficiente sa risolvere.
\end{enumerate}

\section{Ipotesi (i): l'avversario è PPT}

\begin{inbreve}

L'avversario è un programma (una macchina di Turing) che lavora in \textbf{tempo polinomiale} e può usare \textbf{bit casuali}. Si chiama \textbf{PPT}: \emph{Probabilistic Polynomial Time}.

\end{inbreve}

\textbf{Macchina di Turing (TM)} \emph{(definizione non negli appunti del prof)}: il modello matematico di un computer. In pratica: \textbf{un algoritmo}. Dire ``l'avversario è una TM'' vuol dire ``l'avversario è un programma''.

\textbf{Tempo polinomiale}: si conta il numero di passi in funzione della lunghezza dell'input \(n\). Il programma è efficiente se fa \(O(n^c)\) passi per qualche costante \(c \in \mathbb{N}\).

\begin{nota}

\textbf{Cosa vuol dire} \(O(n^c)\): ``al massimo una costante per \(n^c\)''. Esempi: \(n^2\), \(5n^3 + n\), \(n^{100}\) sono polinomiali, cioè efficienti. \(2^n\) \textbf{non} lo è: con \(n = 128\) sono \(2^{128}\) passi, impossibili in pratica. Quindi l'avversario può fare tutto tranne i calcoli esponenziali, come provare tutte le chiavi possibili.

\end{nota}

\textbf{In più}: l'avversario può generare una quantità (polinomiale) di \textbf{bit casuali}, cioè può ``tirare monete'' durante il calcolo.

\[
\text{Adv} \equiv \text{PPT TM} \quad (\text{Probabilistic Polynomial Time Turing Machine})
\]

\section{Ipotesi (ii): esistono problemi difficili}

\begin{inbreve}

Idealmente vorremmo basare la sicurezza su \(P \neq NP\). Nessuno sa farlo; quello che si sa fare è basarla su problemi specifici che si credono difficili, come la \textbf{fattorizzazione}.

\end{inbreve}

\subsection{P e NP}

\begin{inbreve}

\textbf{P} = problemi che si sanno \textbf{risolvere} velocemente. \textbf{NP} = problemi in cui una soluzione, se qualcuno te la dà, si sa \textbf{controllare} velocemente. Ogni problema in P sta anche in NP.

\end{inbreve}

\begin{itemize}
\tightlist
\item
  \(P\) = tutti i problemi risolvibili da una TM in tempo polinomiale.
\item
  \(NP\) = tutti i problemi tali che una soluzione si può \textbf{verificare} da una TM in tempo polinomiale.
\end{itemize}

Esempi del prof di problemi in NP:

\begin{itemize}
\tightlist
\item
  \textbf{FACTORING}: dato \(n = p \cdot q\), trovare \(p\) e \(q\). Se qualcuno ti dà \(p\) e \(q\), controlli subito che \(p \cdot q = n\); trovarli da solo invece è difficile.
\item
  \textbf{3COL}: dato un grafo \(G\), trovare una \textbf{3-colorazione} \(\varphi\) di \(G\). \emph{(Definizione non negli appunti del prof: un grafo è un insieme di punti collegati da linee; una 3-colorazione assegna a ogni punto uno di 3 colori in modo che due punti collegati non abbiano mai lo stesso colore.)} Data una colorazione, controllare che sia giusta è facile (si guarda ogni linea); trovarla è difficile.
\end{itemize}

\begin{esempio}

\textbf{Esempio intuitivo} \emph{(non negli appunti del prof)}: il sudoku. Controllare un sudoku compilato è facile, risolverlo da zero no.

\end{esempio}

\textbf{La domanda P = NP?} È ``tutto ciò che si \textbf{controlla} facilmente si può anche \textbf{trovare} facilmente?''. Nessuno lo sa, ma quasi tutti credono di no (\(P \neq NP\)). P e NP sono \textbf{insiemi di problemi}, non problemi.

\begin{approfondimento}{Approfondimento (non negli appunti del prof): perché le domande sì/no}

Formalmente P e NP parlano di problemi \textbf{sì/no}: ``3233 ha un divisore minore di 60?''. Se la risposta è sì, una \textbf{prova} (es. ``53'') si controlla in fretta. NP = domande in cui la prova del sì si controlla velocemente; P = domande a cui si risponde velocemente da soli. \(P = NP\) vorrebbe dire che non serve mai qualcuno che ti suggerisca la prova: c'è sempre un metodo veloce per arrivarci da solo.

\end{approfondimento}

\subsection{Il teorema ideale}

\[
\textbf{Teorema (ideale)}: \text{ se } P \neq NP, \text{ allora } \Pi \text{ è sicuro.}
\]

\textbf{Perché sarebbe bellissimo}: restano solo due possibilità:

\begin{enumerate}
\def\labelenumi{\arabic{enumi}.}
\tightlist
\item
  \(\Pi\) è sicuro, oppure
\item
  \(P = NP\), che quasi nessuno crede
\end{enumerate}

Rompere \(\Pi\) vorrebbe dire risolvere il problema più famoso dell'informatica.

\textbf{Bad news}: \textbf{nessuno sa dimostrare teoremi di questo tipo.}

\subsection{Quello che si sa fare}

\[
\textbf{Teorema (esempio reale)}: \text{ se FACTORING è difficile, allora } \Pi \text{ è sicuro.}
\]

\textbf{È ancora utile?} Sì, leggendolo al contrario (\textbf{contronominale}): \emph{se} \(\Pi\) \emph{non fosse sicuro, allora FACTORING non sarebbe difficile, cioè esisterebbe una PPT TM che fattorizza.} E siccome da decenni nessuno ha trovato un algoritmo veloce per fattorizzare, è molto improbabile che qualcuno rompa \(\Pi\).

\begin{nota}

\textbf{Contronominale} \emph{(definizione non negli appunti del prof)}: ``se A allora B'' è equivalente a ``se non B allora non A''. Es. ``se piove, la strada è bagnata'' equivale a ``se la strada non è bagnata, non piove''.

\end{nota}

\section{One-way function (OWF)}

\begin{inbreve}

Una \textbf{funzione one-way} è facile da calcolare ma difficile da invertire: dato \(x\) si calcola \(f(x)\) in un attimo, ma dato \(f(x)\) nessun algoritmo efficiente riesce a trovare un input che dia quel risultato. Per il prof è \textbf{``la madre di tutte le ipotesi''}: è l'\textbf{ipotesi minima} per fare crittografia interessante.

\end{inbreve}

\subsection{Spiegazione semplice: la moltiplicazione}

\emph{(Non negli appunti del prof.)}

\textbf{Andata.} Ti do due numeri primi, 61 e 53: moltiplicarli è immediato, \(61 \times 53 = 3233\).

\textbf{Ritorno.} Ti do solo \textbf{3233} e ti chiedo quali due numeri ho moltiplicato. Non c'è un trucco: devi provare a dividere per 2, 3, 5, 7, 11, \ldots{} finché trovi un divisore.

\textbf{Con numeri enormi} (primi da 300 cifre) moltiplicare resta istantaneo, tornare indietro richiederebbe più tempo dell'età dell'universo.

\begin{itemize}
\tightlist
\item
  \(x\) = l'input (i due primi), \(f\) = la funzione (moltiplicarli), \(f(x)\) = il risultato (3233)
\item
  \textbf{invertire} = dato \(f(x)\), trovare \(x\)
\end{itemize}

\begin{nota}

\textbf{Funzione che NON è one-way}: \(f(x) = x + 5\). Da 15 si torna a 10 togliendo 5, istantaneo.

\end{nota}

\textbf{Perché serve}: serve qualcosa che \textbf{tu} sai fare e \textbf{Eve} no. Chi conosce l'input (es. la chiave) calcola tutto facilmente; Eve, che vede solo il risultato, non riesce a tornare indietro.

\subsection{Definizione (del prof)}

Una funzione \textbf{deterministica} \(f : \mathcal{X} \to \mathcal{Y}\) è \textbf{one-way} se per ogni avversario PPT \(A\):

\[
\Pr\big[\text{GAME}^{\text{owf}}_{A,f}(\lambda) = 1\big] \le \text{negl}(\lambda)
\]

cioè la probabilità che \(A\) \textbf{vinca} il gioco qui sotto è \textbf{trascurabile}.

\begin{nota}

\textbf{Cosa sono le cose}

\begin{itemize}
\tightlist
\item
  \textbf{deterministica} = stesso input, stesso output sempre (niente casualità dentro \(f\)).
\item
  \(\mathcal{X}\) = insieme degli input, \(\mathcal{Y}\) = insieme degli output. Tipicamente \(f : \{0,1\}^\lambda \to \{0,1\}^\lambda\): stringhe di \(\lambda\) bit in, stringhe di \(\lambda\) bit out.
\item
  \(\lambda\) = il \textbf{parametro di sicurezza}: la lunghezza dell'input in bit. Più è grande, più il problema è difficile.
\item
  \(\text{GAME}^{\text{owf}}_{A,f}(\lambda)\) = il gioco tra l'avversario \(A\) e il challenger, per la funzione \(f\), con parametro \(\lambda\). Vale 1 se \(A\) vince.
\item
  \(\text{negl}(\lambda)\) = una funzione \textbf{trascurabile} (vedi sotto).
\end{itemize}

\end{nota}

\subsection{Il gioco della one-way function}

\begin{verbatim}
Avversario A                          Challenger C
                                      x ← X      (x scelto a caso)
      ←──────────── y ─────────────   y = f(x)
      ──────────── x' ─────────────→
                                      Output 1 (WIN) se f(x') = y
\end{verbatim}

\begin{enumerate}
\def\labelenumi{\arabic{enumi}.}
\tightlist
\item
  Il challenger sceglie \(x\) \textbf{a caso} in \(\mathcal{X}\) (la freccia \(x \leftarrow \mathcal{X}\) vuol dire ``scelto uniformemente a caso'').
\item
  Calcola \(y = f(x)\) e lo dà ad \(A\).
\item
  \(A\) risponde con un \(x'\).
\item
  \(A\) \textbf{vince} se \(f(x') = y\).
\end{enumerate}

\textbf{Dettagli importanti} \emph{(non negli appunti del prof)}:

\begin{itemize}
\tightlist
\item
  \textbf{Basta una qualunque controimmagine}: \(A\) vince anche se \(x' \neq x\), purché \(f(x') = y\).
\item
  \textbf{L'input è casuale}: la difficoltà è richiesta \textbf{in media}, su input casuali, non solo per qualche input particolare.
\end{itemize}

\subsection{Esempi di OWF (candidate)}

Nessuno ha mai \textbf{dimostrato} che le OWF esistono: sono un'\textbf{ipotesi}. Questi sono i candidati del prof, cioè problemi che si credono difficili:

\begin{itemize}
\tightlist
\item
  \textbf{FACTORING}: \(f(p, q) = p \cdot q\); invertire = fattorizzare.
\item
  \textbf{DISCRETE LOG} (logaritmo discreto): \(f(x) = g^x \bmod p\); invertire = trovare l'esponente \(x\) dato \(g^x \bmod p\).
\item
  \textbf{LATTICES} (reticoli) e \textbf{DECODING RANDOM CODES} (decodificare codici casuali): problemi di algebra lineare ``rumorosa''. \emph{(Non negli appunti del prof: sono importanti perché si crede che resistano anche ai computer quantistici.)}
\end{itemize}

\subsection{OWF e P vs NP}

\emph{(Non negli appunti del prof: approfondimento per capire perché le OWF sono un'ipotesi più forte di P ≠ NP.)}

\begin{inbreve}

Se le OWF esistono, allora \(P \neq NP\). Il viceversa non è garantito: anche se \(P \neq NP\) le OWF potrebbero non esistere.

\end{inbreve}

\begin{approfondimento}{Dimostrazione: OWF ⇒ P ≠ NP}

Si dimostra al contrario: \textbf{se P = NP, allora nessuna funzione è one-way}.

P e NP parlano di domande sì/no, ma a noi serve \textbf{trovare} \(x\). Si fa come nel gioco delle 20 domande: si scopre \(x\) \textbf{un bit alla volta}, con domande del tipo ``esiste un \(x'\) che inizia con \(p\) e tale che \(f(x') = y\)?''.

Questa domanda sta in NP: se qualcuno propone \(x'\), si controlla velocemente (si guarda come inizia e si calcola \(f(x')\)). Se \(P = NP\), c'è un algoritmo veloce che risponde.

\textbf{L'invertitore}: parte dal prefisso vuoto; a ogni passo chiede ``esiste una controimmagine che continua con 0?''; se sì aggiunge 0, altrimenti aggiunge 1. Resta sempre vero che esiste una controimmagine con quel prefisso, quindi dopo \(\lambda\) passi ha una controimmagine completa.

\textbf{Esempio} con \(x = 101\): inizia con 0? No → 1. Inizia con 10? Sì → 10. Inizia con 100? No → 101 ✓.

Sono \(\lambda\) domande veloci, quindi tempo polinomiale, e successo con probabilità 1: \(f\) non è one-way. ∎

\end{approfondimento}

\textbf{Perché P ≠ NP non basta}: \(P \neq NP\) garantisce solo che esiste \textbf{almeno qualche} caso difficile (caso peggiore). In crittografia la chiave è scelta \textbf{a caso}, quindi serve che il problema sia difficile \textbf{quasi sempre} (caso medio). È come un produttore che garantisce che almeno uno su 1000 lucchetti è impossibile da scassinare: se ne prendi uno a caso, quasi sicuramente è facile.

\section{Funzioni trascurabili (negligible)}

\begin{inbreve}

``piccolo'' in crittografia si definisce in modo \textbf{asintotico}: una funzione è \textbf{trascurabile} se diventa più piccola dell'inverso di \textbf{qualunque} polinomio. È il modo formale di dire ``probabilità così piccola che si può ignorare''.

\end{inbreve}

\subsection{Definizione (del prof)}

Una funzione \(\varepsilon(\lambda)\) è \textbf{trascurabile} (\emph{negligible}) in \(\lambda\) se

\[
\varepsilon(\lambda) = O\Big(\frac{1}{p(\lambda)}\Big) \quad \text{per } \textbf{ogni} \text{ polinomio } p(\lambda) = \text{poly}(\lambda)
\]

Si scrive \(\varepsilon(\lambda) = \text{negl}(\lambda)\).

\textbf{Come leggerla} \emph{(non negli appunti del prof)}: per ogni polinomio, anche enorme come \(\lambda^{100}\), da un certo punto in poi \(\varepsilon(\lambda)\) sta sotto \(1/p(\lambda)\) e ci resta. Basta \textbf{un solo} polinomio per cui non succede e la funzione non è trascurabile.

\subsection{Esempi}

\textbf{Esempio del prof}: \(2^{-\lambda}\) è trascurabile. Confronto con \(1/\lambda^{10}\):

\begin{longtable}[]{@{}llll@{}}
\toprule\noalign{}
λ & 2\^{}(−λ) & 1/λ\^{}10 & Chi è più piccolo? \\
\midrule\noalign{}
\endhead
\bottomrule\noalign{}
\endlastfoot
10 & ≈ 0,001 & 0,0000000001 & 1/λ\^{}10 \\
50 & ≈ 9 · 10\^{}(−16) & ≈ 1 · 10\^{}(−17) & 1/λ\^{}10 \\
100 & ≈ 8 · 10\^{}(−31) & 1 · 10\^{}(−20) & \textbf{2\^{}(−λ)} \\
200 & ≈ 6 · 10\^{}(−61) & ≈ 1 · 10\^{}(−23) & \textbf{2\^{}(−λ)} \\
\end{longtable}

All'inizio \(2^{-\lambda}\) è più grande, ma da \(\lambda = 59\) in poi passa sotto e non risale più: si dimezza a ogni passo, mentre \(1/\lambda^{10}\) scende sempre più lentamente. Con qualunque altro polinomio succede lo stesso, solo più tardi.

\textbf{Controesempio} \emph{(non negli appunti del prof)}: \(1/\lambda^{10}\) \textbf{non} è trascurabile. Con il polinomio \(\lambda^{11}\) dovrebbe finire sotto \(1/\lambda^{11}\), ma \(1/\lambda^{10}\) è sempre più grande (es. con \(\lambda = 10\): \(10^{-10} > 10^{-11}\)).

\textbf{Perché conta}: un attacco che riesce con probabilità \(1/\lambda^{10}\) si può ripetere \(\lambda^{10}\) volte (tempo ancora polinomiale) e alla fine riesce. Con probabilità \(2^{-\lambda}\), anche ripetendolo \(\lambda^{10}\) volte si arriva al massimo a \(\lambda^{10} \cdot 2^{-\lambda}\), ancora minuscolo.

\subsection{Perché è comoda (proprietà del prof)}

\[
\text{negl}(\lambda) + \text{negl}(\lambda) = \text{negl}(\lambda) \qquad \text{poly}(\lambda) \cdot \text{negl}(\lambda) = \text{negl}(\lambda)
\]

\begin{approfondimento}{Dimostrazione delle due proprietà (non negli appunti del prof)}

\textbf{Idea}: per battere \(1/\lambda^c\) si usa la definizione con un esponente \textbf{più grande} di \(c\), così c'è margine per assorbire il fattore 2 o il polinomio.

\textbf{Somma.} Fissato \(c\), si applica la definizione con \(c+1\): da un certo punto \(\varepsilon_1, \varepsilon_2 < 1/\lambda^{c+1}\). Allora per \(\lambda \ge 2\): \(\varepsilon_1 + \varepsilon_2 < 2/\lambda^{c+1} = (2/\lambda) \cdot 1/\lambda^c \le 1/\lambda^c\). ✓

\textbf{Prodotto.} Un polinomio da un certo punto sta sotto \(\lambda^d\) per qualche \(d\). Fissato \(c\), si applica la definizione con \(c + d\): \(\varepsilon < 1/\lambda^{c+d}\), quindi \(p(\lambda) \cdot \varepsilon < \lambda^d / \lambda^{c+d} = 1/\lambda^c\). ∎

\end{approfondimento}

\section{Conclusione: i mondi computazionali}

\begin{inbreve}

A seconda delle ipotesi che si fanno, si ottengono ``mondi'' diversi, cioè diverse cose che la crittografia riesce a costruire. Con le sole OWF si costruisce tutta la crittografia \textbf{simmetrica}; per quella \textbf{a chiave pubblica} sembra servire qualcosa in più.

\end{inbreve}

\begin{longtable}[]{@{}
  >{\raggedright\arraybackslash}p{(\columnwidth - 4\tabcolsep) * \real{0.3333}}
  >{\raggedright\arraybackslash}p{(\columnwidth - 4\tabcolsep) * \real{0.3333}}
  >{\raggedright\arraybackslash}p{(\columnwidth - 4\tabcolsep) * \real{0.3333}}@{}}
\toprule\noalign{}
\begin{minipage}[b]{\linewidth}\raggedright
Mondo
\end{minipage} & \begin{minipage}[b]{\linewidth}\raggedright
Ipotesi
\end{minipage} & \begin{minipage}[b]{\linewidth}\raggedright
Cosa si può costruire
\end{minipage} \\
\midrule\noalign{}
\endhead
\bottomrule\noalign{}
\endlastfoot
\textbf{Minicrypt} & Esistono le OWF & SKE, MAC, PRNG \\
\textbf{Cryptomania} & Ipotesi più forti delle OWF & PKE, scambio di chiavi (TLS) \\
\end{longtable}

\begin{nota}

\textbf{Le sigle} \emph{(definizioni non negli appunti del prof)}

\begin{itemize}
\tightlist
\item
  \textbf{SKE} (\emph{Symmetric Key Encryption}): cifratura a chiave simmetrica, Alice e Bob hanno la stessa chiave (come l'OTP).
\item
  \textbf{MAC}: codici di autenticazione dei messaggi (lezioni 2-3).
\item
  \textbf{PRNG} (\emph{Pseudo-Random Number Generator}): generatori di bit pseudocasuali, cioè che \emph{sembrano} casuali (vedi Lezione 5).
\item
  \textbf{PKE} (\emph{Public Key Encryption}): cifratura a chiave pubblica, una chiave pubblica per cifrare e una privata per decifrare (es. RSA).
\item
  \textbf{Scambio di chiavi} (\emph{key exchange}): due persone che non si sono mai viste concordano una chiave segreta parlando su un canale pubblico. È quello che fa \textbf{TLS} ogni volta che apri un sito in HTTPS.
\end{itemize}

\end{nota}

\textbf{La freccia sbarrata} nel disegno del prof: dalle sole OWF \textbf{non} si sa ottenere la crittografia a chiave pubblica. Minicrypt non porta automaticamente a Cryptomania.

\begin{approfondimento}{I cinque mondi di Impagliazzo (dalle dispense, non negli appunti del
prof)}

Russell Impagliazzo (1995) ha descritto cinque mondi possibili. Non sappiamo in quale viviamo; si spera in uno degli ultimi due.

\begin{itemize}
\tightlist
\item
  \textbf{Algorithmica}: \(P = NP\), niente problemi difficili, niente crittografia computazionale.
\item
  \textbf{Heuristica}: \(P \neq NP\), ma i problemi difficili sono facili \textbf{in media}: inutile per le chiavi casuali.
\item
  \textbf{Pessiland}: esistono problemi difficili in media, ma \textbf{non} le OWF. Il mondo peggiore.
\item
  \textbf{Minicrypt}: esistono le OWF, ma non la crittografia a chiave pubblica.
\item
  \textbf{Cryptomania}: esistono OWF \textbf{e} crittografia a chiave pubblica.
\end{itemize}

\end{approfondimento}

\section{Riepilogo}

\begin{itemize}
\tightlist
\item
  \textbf{Lemma della collision probability}: si dimostra scrivendo la SD come prodotto scalare e usando \textbf{Cauchy-Schwarz}; l'eccesso di collisioni controlla la distanza dall'uniforme
\item
  \textbf{Sicurezza computazionale}: due ipotesi in più, avversario efficiente e problemi difficili
\item
  \textbf{Avversario = PPT TM}: tempo polinomiale \(O(n^c)\) e bit casuali
\item
  \textbf{P} = problemi che si risolvono velocemente; \textbf{NP} = problemi in cui una soluzione si controlla velocemente (es. FACTORING, 3COL)
\item
  Ideale: ``se P ≠ NP allora Π sicuro'', ma nessuno sa farlo. Reale: ``se FACTORING è difficile allora Π sicuro'', e al contrario: rompere Π ⇒ fattorizzare
\item
  \textbf{OWF}: facile da calcolare, difficile da invertire; nel gioco l'avversario vince con probabilità \(\le \text{negl}(\lambda)\); è l'ipotesi minima della crittografia
\item
  \textbf{Trascurabile}: \(O(1/p(\lambda))\) per ogni polinomio; es. \(2^{-\lambda}\); negl + negl = negl, poly · negl = negl
\item
  Candidati OWF: FACTORING, DISCRETE LOG, LATTICES, DECODING RANDOM CODES
\item
  \textbf{Minicrypt} (OWF): SKE, MAC, PRNG. \textbf{Cryptomania}: PKE, scambio di chiavi (TLS)
\end{itemize}
